\documentclass[a4paper, 11pt]{article}

% Pacchetti fondamentali per codifica, lingua e margini
\usepackage[utf8]{inputenc}
\usepackage[T1]{fontenc}
\usepackage[italian]{babel}
\usepackage[margin=2.5cm]{geometry}

% Font simile a quello del PDF originale (Times/Liberation Serif)
\usepackage{tgtermes}

% Pacchetti per immagini, colori e tabelle
\usepackage{graphicx}
\usepackage[table]{xcolor}
\usepackage{array}
\usepackage{tabularx}

% Pacchetti per la formattazione dei paragrafi, titoli e indice
\usepackage[parfill]{parskip}
\usepackage{titlesec}
\usepackage{tocloft}
\usepackage[hidelinks]{hyperref}

% Definizione del colore rosso UniPD
\definecolor{unipdred}{RGB}{155, 0, 20}

% Formattazione personalizzata per i sottoparagrafi (2.1, 2.2, ecc.)
\titleformat{\subsection}
  {\normalfont\normalsize\bfseries}{\thesubsection}{0.5em}{}
\titlespacing*{\subsection}{0pt}{1.5ex plus 0.5ex minus .2ex}{0.2ex}

% Configurazione dei puntini nell'indice anche per le sezioni principali
\renewcommand{\cftsecleader}{\cftdotfill{\cftdotsep}}
\renewcommand{\cftsecfont}{\normalfont}
\renewcommand{\cftsecpagefont}{\normalfont}

\begin{document}

%INTESTAZIONE COHI LOGHI
\noindent
\begin{minipage}[c]{0.48\textwidth}
    \begin{tabular}{@{}m{2.6cm} m{5cm}@{}}
        % aggiungere immagine
        \includegraphics[width=2.4cm]{logo_unipd.png} & 
        \textcolor{unipdred}{
            Università di Padova \newline
            Corso di \newline
            Ingegneria del Software \newline
            A.A. 2026/2027 
        }
    \end{tabular}
\end{minipage}
\hfill
\begin{minipage}[c]{0.48\textwidth}
    \begin{tabular}{@{}m{2.4cm} m{5.5cm}@{}}
        % aggiungere immagine
        \includegraphics[width=2.2cm]{logo_gruppo.jpeg} & 
\begin{figure}
            \centering
            \includegraphics[width=0.5\linewidth]{logo_unipd.png}
            \caption{Enter Caption}
\begin{figure}
                \centering
                \includegraphics[width=0.5\linewidth]{logo_gruppo.jpeg}
                \caption{Enter Caption}
                \label{fig:placeholder}
            \end{figure}
                        \label{fig:placeholder}
        \end{figure}
                Gruppo 18 \newline
        The Lord of the Strings \newline
        thelordofthestrings@gmail.com
    \end{tabular}
\end{minipage}

\vspace{3cm}

%TITOLO DEL DOCUMENTO
{\LARGE Verbale interno della riunione in data 2000/01/01\par}

\vspace{1.8cm}

%INDICE

\phantomsection
\addcontentsline{toc}{section}{\protect\numberline{1}Informazioni generali}

\tableofcontents

\vspace{1cm}

%DATI E STORICO MODIFICHE
\textbf{Redatto da: Pinco Palla}

\vspace{0.4cm}

\textbf{Storico modifiche}

\vspace{-0.2cm}
\noindent
\renewcommand{\arraystretch}{1.3}
\begin{tabularx}{\textwidth}{|X|X|X|X|}
    \hline
    \rowcolor{gray!35}
    Versione & Data & Descrizione & Autore \\
    \hline
    1.0.0 & 2000/01/01 & Aggiungi & Pinco Palla \\
    \hline
\end{tabularx}

%SECONDA PAGINA
\newpage

% Imposta il contatore alla sezione 2 per generare 2.1, 2.2, 2.3.

\setcounter{section}{2}
\addcontentsline{toc}{section}{\protect\numberline{2}}

\subsection{Durata e partecipanti}
Ora: 00:00 – 00:00 \\
Partecipanti:
\begin{itemize}\setlength\itemsep{0em}
    \item Giovanni Biba
    \item Giacomo De Lucchi
    \item Costanza Gyulai
    \item Michele Moro
    \item Stefano Sironi
    \item Angelo Tuoni
    \item Yi Ming Ye
\end{itemize}
Mezzo: Discord

\vspace{0.4cm}

\subsection{Sintesi}
Breve resoconto di cosa è stato detto e deciso.

\vspace{0.4cm}

\subsection{Obbiettivi fissati}
Breve descrizione di cosa è stato deciso e stabilito

\end{document}